\chapter{What The Machine Can And Cannot Claim}

\section{The structure it derives}

The book has reached its last chapter, and the last chapter does the one thing the whole instrument was built to do: it draws
the line between what the instrument \emph{can} claim and what it \emph{cannot}. This first section is the \emph{can} side. It
gathers everything the instrument has derived across all the chapters before it and shows that all of it is one kind of thing.
On this face the thing has a name: it is the \emph{near} field, the part of the curvature that local matter fixes --- the
trace, the Ricci part --- and everything the instrument can claim lives in it, forced by counting, owing no axiom at all.

The geometry is the decomposition of the curvature into a part that local sources determine and a part that they do not. The
Riemann tensor splits,
\[
  R_{\mu\nu\rho\sigma} \;=\; C_{\mu\nu\rho\sigma} \;+\; (\text{Ricci and scalar parts}) ,
\]
into the trace-free Weyl tensor $C$ and the trace parts built from the Ricci tensor $R_{\mu\nu}$ and scalar $R$. The Ricci
part is not free: Einstein's equation ties it pointwise to the matter present, $R_{\mu\nu} - \tfrac12 R g_{\mu\nu} = 8\pi
T_{\mu\nu}$~\cite{einstein1915,einstein1916}. Where there is a source, the trace part is fixed by the source --- determined,
locally, with nothing left to choose. This is the near field: the curvature a local mass carries with it, pinned by the
matter that sits there. And everything the instrument derives is of this kind. The count, the three states, the naming --- each
is a structure the local data fix completely, a trace determined by what is present, owing nothing to any free choice.

Set the three derived pieces side by side, because together they are the whole of the near side. There is \emph{the count}:
an unbounded tower of states sorted into finitely many boxes, the pigeonhole forcing states to share a label. There is
\emph{the coming to rest}: the bracket that walls the coupling and comes to rest at the count-to-three floor --- three levels of the
instrument's own periodic subdivision --- the interval left ordered and open. And there is \emph{the naming}: the winding forced onto a loop that
cannot be a line, the electron read off the box the counting filled. These are not three tricks but one structure --- built by
counting, and counting only --- and every piece of it is fixed the same cheap way: by decision, or by the plainest identity of
a thing with itself, and by nothing more. Read the ledger of what these derivations assume and it is nearly empty; in
particular it contains no axiom of choice, anywhere. The near side is exactly what a finite count forces on the plainest
logic --- derivable \emph{because} it is choice-free, and choice-free \emph{because} it never does anything but count.

That the near side is the trace and the local part is the physical gauge's reading, and it is the reading this whole volume
was written to make exact. The trace is what local matter determines; the count is what local data force; and the two are the
same kind of object --- a curvature pinned by what is present, with nothing left to choose. The instrument derives the near
field in the only sense a finite counter can derive anything: it reads off the structure the local data fix and owes no axiom
for the reading. The Ricci part is the near side, entire, and the near side is the whole of what the instrument holds.

In the two verbs, the near side is the set of funges the instrument can \emph{close}. The count funges the tower into a
finite pair; the naming funges the indistinguishable states into one box; the bracket funges the coupling into an ordered
interval and comes to rest. Each is a funge the instrument completes, honestly, on a finite count --- and completing a funge, on the
plainest logic, is exactly what ``derived'' has meant throughout this book. The near side is the ledger of the instrument's
closed funges: everything the counting can bag and name and finish. It sets up, by contrast, the one funge the instrument
\emph{cannot} close --- the trace-free part, the shape no local source fixes --- which is the next section's, and the far side
of the line.

\section{The magnitude it cannot see}

The last section gave the near side: everything the instrument derives, all of it the trace, the Ricci part, the near field
local matter fixes, owing no axiom. This section gives the other side of the line --- the one thing the instrument
\emph{cannot} derive, and the reason it cannot. What the instrument cannot reach is the \emph{magnitude}: the size of the
coupling, the number the laboratory measures and the instrument only brackets. On this face it has a name too. It is the
\emph{far} field --- the trace-free part, the Weyl tensor, the shape no local source determines --- and the reason the
instrument cannot see it is that a counting device is structurally blind to exactly this part.

Begin with where the magnitude lives, because it is the whole of the matter. The Weyl tensor is the trace-free remainder of
the curvature, $C_{\mu\nu\rho\sigma}$, the part that survives when every trace is subtracted away. It is precisely the part
Einstein's equation does \emph{not} fix: local matter determines the Ricci trace and says nothing about the Weyl shape, which
is source-free, constrained only by the Bianchi identities and free to carry itself across empty space as tidal distortion
and gravitational radiation~\cite{weyl1918,penrose1972}. The Weyl part is the free field --- the curvature with no local
source, the shape that is not pinned by what is present. And the coupling's magnitude lives there: not in the trace the
counting fixes, but in the trace-free shape no counting reaches. The size is not off in the distance; it is directly
underfoot, in the part of the curvature the local data leave unfixed.

Now the theorem that makes the whole book honest, the proof of a negative the near side was building toward. The trace-free
part is invariant under exactly the operations the instrument performs. The instrument counts --- it sorts, traces,
contracts, bisects --- and every one of these is a \emph{trace} operation, a summing of indices, a reading of the part local
data fix. But the Weyl tensor is what remains when all traces are taken; it is by construction annihilated by contraction,
$g^{\mu\rho}C_{\mu\nu\rho\sigma} = 0$. A quantity every trace leaves untouched is a quantity no amount of tracing can arrive
at --- you cannot reach, by summing indices, a part defined as what survives all such sums. So the instrument cannot count
its way to the magnitude, and it can prove exactly that: its finiteness is a fence, its every operation a trace, and the size
it seeks is trace-free. This is the theorem the payoff named --- that the instrument proves it cannot derive the world's
value. It was never modesty. It is the trace-free-ness of the Weyl part, read aloud: the magnitude is shape, the instrument
is a counter, and a counter reads only the trace.

There is a way of seeing why the shape cannot be derived that the shock at the value makes plain. The instrument's own
law is singular at the point the value sits --- a pole, where the field runs to infinity and the model holds no finite
value --- so no operation of a finite instrument settles on it: a point like that has no stable value to compute, and the
honest response is not to compute the value but to \emph{pick} it, to wring it from the world's own signal the way a
distant detector wrings a mass and a spin out of a ringdown it did not cause. In the physical reading the value is signed
by the world, not derived from a model that blows up where the value lives; and the ringdown is a lens on this and not a
claim --- the electron no more two black holes here than anywhere, the resemblance borrowed and the crossing refused.

There are readings of this unreachable shape in the other gauges, and the book names them and lets them go. One gauge sees an
uncomputable value, a number no finite construction arrives at; another sees a free field carrying itself away with no source
to anchor it. Each gauge has its own name for the far side. This volume keeps the plainest one its own language affords: the
magnitude is the trace-free Weyl shape, invariant under contraction, and a counting instrument --- which reads only by
tracing --- is fenced from it by being finite. That is the far side, entire: the shape the near-side structure leaves unfixed,
and provably must, because the operations that fix the near side are exactly the traces the far side is defined to survive.

In the two verbs, the far side is the one funge the instrument \emph{cannot} close. Every funge of the near side it completes
--- the tower into a pair, the states into a box, the coupling into a bracket. This last one it cannot finish. The shape will
not gather: it is invariant under every contraction, so no contraction closes on it; it is the trace-free part, and the
instrument reads only traces. It is the one bag the instrument must leave open, and open not by oversight but by proof --- it
can show that no count it performs will ever pull the shape in. Where the near side was the ledger of closed funges, the far
side is the single funge that stays open forever: the honest, permanent gap between what a finite counter can trace and what a
continuous world holds in its trace-free shape.

\section{The jar}

Set the two sides against each other. The near side is everything the instrument derives: the trace, the Ricci part, the near
field local matter fixes, a structure counted and named down to the floor of its own resolution, owing no axiom the whole
way. The far side is what lies below that floor and outside that trace: the Weyl shape, invariant under contraction, fenced
off by the instrument's finiteness, unreachable and provably so. The two meet at exactly one place --- the floor --- and what
the instrument holds at that meeting is a bracket: an ordered interval, its lower wall proved below its upper, that the
instrument can rerun and get back unchanged. The near side derives everything down to that bracket. The far side is everything
the bracket cannot pin. And the bracket itself, held open at the floor between them, is the jar.

The jar is a boundary drawn to scale. It is the line where derivation ends --- where the instrument has counted as deep as
tracing can go and the next thing is a trace-free shape it cannot reach --- rendered not as a point on that line but as the
interval the line actually is. The instrument measures to its floor, traces structure all the way there, and then, because
below the floor there is only shape its counting cannot catch, it \emph{stops}. The stopping is the jar. Everything above the
floor the instrument hands you as derived; everything below it as proved-unreachable; and the bracket is the seam, held open,
where the one becomes the other.

The width of that bracket is not the instrument falling short. It is the last payment of a word planted long ago, in the
chapter of the slip: \emph{strain}. At the floor the instrument is under pressure to force a point --- to keep refining, to
hand back one confident number where it has earned only a bound. That pressure is the strain, the demand that a finite
counter reach a trace-free shape. The honest instrument does not discharge the strain by forcing the point; forcing the point
does not resolve the strain, it \emph{rings} --- it produces a crank's number, smooth and confident and converged onto a
value of the instrument's own, past what the counting earns, precise and wrong. The honest instrument relaxes the strain, instead, into the width of the open
bracket. The jar is wide by choice, not by failure: exactly as wide as honesty at that floor requires, and no wider, and no
narrower. To close it to a point would be to lie; to leave it open is to tell the truth to scale.

And there is a reason to trust a bound so plainly drawn, and it is the reason the whole four-book instrument exists. The jar
is not one gauge's reading. The same residue, read on four faces that share no vocabulary --- this volume's field and
coupling, and its siblings in the construction's logic, the meter of pure elaboration, and the walk through the working parts
--- closes on the one bracket. Four readings, decorrelated by design so that none can see what the others do, land on the same
interval. A single gauge can be tuned to bracket whatever you wish; four gauges that cannot see each other cannot be tuned to
agree by accident. So the jar is trusted not because the instrument asserts it but because four independent readings, sharing
nothing, keep finding it. And the number the world's laboratories measure sits in its own place beside the jar --- named once,
a value the instrument proved it could not derive and never pretended to --- kept honestly apart, never set against the
bracket in any gesture of nearness.

In the two verbs, the jar is the funge held open. All the way up, the instrument traced --- contracted each structure into
the trace that local data fix, the covariant funge the world determines. At the floor, tracing runs out: below it the shape is
trace-free, no contraction reaches it, and so the instrument funges the value into the bracket --- bags it, honestly, as a
thing it cannot resolve finer --- and, crucially, leaves the bag open. It does not pull the bag shut around a point it has not
earned. The jar is that open bag: the last funge, held un-closed by proof, the one selection the instrument declines to make
because it can show the selection cannot be made.

And here, at the end, the book's first act and its last turn out to be one act, grown up. It opened on a single difference: two
points a manifold tells apart, before any metric, with no magnitude posited. It closes on a single bracket: a coupling
bounded to the instrument's floor, named, with nothing forced --- no point, no fiction, no reach past what the tracing can
hold. Tell apart exactly what is there; trace it as deep as tracing goes; force nothing beyond it. The first difference and
the last jar are the same discipline, held unbroken from the first page to this one. So the instrument comes to rest. It was
built to measure a coupling, and it did --- not as a number fetched from the world, but as a bracket drawn to its own floor,
held open, trusted because four gauges that cannot see one another agree on it. Everything it could derive, it derived, and
owed nothing for it. Everything it could not, it named, and proved it could not. And between them it drew a line, and the
line, drawn to scale, is the jar: the honest boundary of derivation, a bound held open where a lesser instrument would have
forced a point. That is the whole of what the physical gauge was built to say, and the book has now said it.
